\subsection{Controlled ablation on a locked representative set}

We evaluated the agent on a locked 5-case subset (cases 01, 02, 09, 14, and 19), chosen to cover atomic H, NNH, OH, and acetaldehyde adsorption across increasingly difficult alloy surfaces. For each case, we executed five runtime variants under a fixed MACE-small CPU backend: \textit{Full}, \textit{$-$Slip}, \textit{$-$Forbid}, \textit{$-$Term}, and \textit{1-Shot}. Here \textit{$-$Slip} removes slip-conditioned guidance, \textit{$-$Forbid} keeps slip detection but removes the explicit ``do not retry this site'' instruction, \textit{$-$Term} disables convergence-based early stopping, and \textit{1-Shot} reduces the total budget to a single attempt.

% ============================================================
% AdsMind Paper-Ready Tables
% Regenerated after completed 4-backend ablation matrices
% Source artifacts:
%   - research/results/gemini_ablation_v1/ablation_summary.csv
%   - research/results/xai_ablation_v2/ablation_summary.csv
%   - research/results/openai_gpt54_ablation_v1/ablation_summary.csv
%   - research/results/anthropic_sonnet46_ablation_v1/ablation_summary.csv
% ============================================================

% Table 1: Ablation Energy Matrix (Gemini)
\begin{table}[t]
\centering
\caption{Gemini ablation study on the five locked representative cases. Lower (more negative) adsorption energy is better. Values are rounded to 3 decimals; ties at the displayed precision are co-bolded.}
\label{tab:ablation-gemini}
\begin{tabular}{lrrrrr}
\toprule
Case & Full & $-$Slip & $-$Forbid & $-$Term & 1-Shot \\
\midrule
01 & \textbf{$-$3.632} & \textbf{$-$3.632} & \textbf{$-$3.632} & \textbf{$-$3.632} & $-$3.631 \\
02 & \textbf{$-$4.766} & \textbf{$-$4.766} & \textbf{$-$4.766} & \textbf{$-$4.766} & \textbf{$-$4.766} \\
09 & \textbf{$-$1.974} & \textbf{$-$1.974} & \textbf{$-$1.974} & \textbf{$-$1.974} & $-$1.801 \\
14 & \textbf{$-$3.617} & \textbf{$-$3.617} & \textbf{$-$3.617} & \textbf{$-$3.617} & \textbf{$-$3.617} \\
19 & $-$3.841 & $-$3.841 & \textbf{$-$4.042} & $-$3.594 & $-$3.133 \\
\bottomrule
\end{tabular}
\end{table}

% Table 2: Ablation Energy Matrix (Grok-4)
\begin{table}[t]
\centering
\caption{Grok-4 ablation study on the same five representative cases. Lower (more negative) adsorption energy is better. Values are rounded to 3 decimals; ties at the displayed precision are co-bolded.}
\label{tab:ablation-grok}
\begin{tabular}{lrrrrr}
\toprule
Case & Full & $-$Slip & $-$Forbid & $-$Term & 1-Shot \\
\midrule
01 & \textbf{$-$3.632} & \textbf{$-$3.632} & \textbf{$-$3.632} & \textbf{$-$3.632} & $-$3.631 \\
02 & \textbf{$-$4.766} & \textbf{$-$4.766} & $-$4.760 & \textbf{$-$4.766} & $-$4.150 \\
09 & \textbf{$-$1.974} & \textbf{$-$1.974} & \textbf{$-$1.974} & \textbf{$-$1.974} & \textbf{$-$1.974} \\
14 & \textbf{$-$3.617} & \textbf{$-$3.617} & \textbf{$-$3.617} & \textbf{$-$3.617} & $-$3.245 \\
19 & \textbf{$-$4.045} & $-$3.594 & $-$3.594 & $-$4.027 & $-$3.594 \\
\bottomrule
\end{tabular}
\end{table}

% Table 3: Ablation Energy Matrix (GPT-5.4)
\begin{table}[t]
\centering
\caption{GPT-5.4 ablation study on the same five representative cases. Lower (more negative) adsorption energy is better. Values are rounded to 3 decimals; ties at the displayed precision are co-bolded.}
\label{tab:ablation-gpt54}
\begin{tabular}{lrrrrr}
\toprule
Case & Full & $-$Slip & $-$Forbid & $-$Term & 1-Shot \\
\midrule
01 & $-$3.627 & $-$3.627 & $-$3.627 & \textbf{$-$3.630} & $-$3.557 \\
02 & \textbf{$-$4.769} & $-$4.764 & \textbf{$-$4.769} & \textbf{$-$4.769} & $-$4.343 \\
09 & \textbf{$-$1.991} & \textbf{$-$1.991} & \textbf{$-$1.991} & \textbf{$-$1.991} & \textbf{$-$1.991} \\
14 & \textbf{$-$3.591} & $-$3.431 & \textbf{$-$3.591} & \textbf{$-$3.591} & $-$3.431 \\
19 & $-$4.013 & $-$4.100 & \textbf{$-$4.149} & $-$4.013 & $-$4.013 \\
\bottomrule
\end{tabular}
\end{table}

% Table 4: Ablation Energy Matrix (Claude Sonnet 4.6)
\begin{table}[t]
\centering
\caption{Claude~Sonnet~4.6 ablation study on the same five representative cases. Lower (more negative) adsorption energy is better. Values are rounded to 3 decimals; ties at the displayed precision are co-bolded.}
\label{tab:ablation-claude}
\begin{tabular}{lrrrrr}
\toprule
Case & Full & $-$Slip & $-$Forbid & $-$Term & 1-Shot \\
\midrule
01 & \textbf{$-$3.630} & \textbf{$-$3.630} & $-$3.624 & \textbf{$-$3.630} & $-$3.557 \\
02 & \textbf{$-$4.769} & \textbf{$-$4.769} & \textbf{$-$4.769} & \textbf{$-$4.769} & $-$4.343 \\
09 & \textbf{$-$1.991} & \textbf{$-$1.991} & \textbf{$-$1.991} & \textbf{$-$1.991} & $-$1.804 \\
14 & $-$3.580 & \textbf{$-$3.591} & \textbf{$-$3.591} & \textbf{$-$3.591} & $-$3.453 \\
19 & $-$4.013 & \textbf{$-$4.149} & \textbf{$-$4.149} & $-$4.013 & $-$4.013 \\
\bottomrule
\end{tabular}
\end{table}

% Table 5: Cross-LLM robustness summary
\begin{table}[t]
\centering
\caption{Cross-LLM robustness across the completed 5$\times$5 ablation matrix. The four-backend range is max--min energy across Gemini, Grok-4, GPT-5.4, and Claude for the same case and runtime variant.}
\label{tab:cross-llm-summary}
\begin{tabular}{lrrr}
\toprule
Variant & Mean 4-backend range (eV) & $\leq$0.01 eV & Largest range (eV) \\
\midrule
Full & 0.053 & 2/5 & 0.203 \\
$-$Slip & 0.153 & 2/5 & 0.555 \\
$-$Forbid & 0.123 & 2/5 & 0.555 \\
$-$Term & 0.096 & 2/5 & 0.433 \\
1-Shot & 0.426 & 0/5 & 0.880 \\
\midrule
Overall & 0.170 & 8/25 & 0.880 \\
\bottomrule
\end{tabular}
\end{table}

% Table 5: Key evaluation metrics
\begin{table}[t]
\centering
\caption{Key evaluation metrics across all four backends. $^\dagger$GPT-5.4 and Claude one-shot slip-rate denominators reflect valid (non-dissociated) results only.}
\label{tab:key-evaluation-metrics}
\begin{tabular}{p{2.4cm}p{1.1cm}p{2.1cm}p{2.1cm}p{2.1cm}p{2.1cm}p{1.0cm}}
\toprule
Metric & Target & Gemini & Grok-4 & GPT-5.4 & Claude & Met \\
\midrule
1-shot slip rate & $< 20\%$ & 60\% (12/20) & 60\% (12/20) & 63\% (12/19)$^\dagger$ & 67\% (12/18)$^\dagger$ & No \\
Friedman $p$ & $< 0.05$ & 0.30 & 0.018 & 0.062 & 0.010 & 2/4 \\
H$_1$ support & $\geq$75\% & 1/4 (25\%) & 3/4 (75\%) & 3/4 (75\%) & 3/4 (75\%) & 3/4 \\
Backend conv. & Lower & \multicolumn{4}{p{8.4cm}}{Full mean 4-backend range 0.053~eV vs.\ 0.426~eV 1-shot; 8.0$\times$ reduction.} & Yes \\
\bottomrule
\end{tabular}
\end{table}

% Table 6: Explicit H1 hypothesis test
\begin{table}[t]
\centering
\caption{Explicit test of H$_1$ on the locked ablation subset using $\Delta E = E_{\mathrm{full}} - E_{\mathrm{1shot}}$. Negative values mean the full closed-loop agent found a lower energy; hits require $\Delta E < -0.05$~eV. H$_1$ uses only intermetallic surfaces; case~09 is shown for completeness but excluded from the support rate.}
\label{tab:hypothesis-test}
\begin{tabular}{llrrrr}
\toprule
Case & Family & Gemini $\Delta E$/hit & Grok-4 $\Delta E$/hit & GPT-5.4 $\Delta E$/hit & Claude $\Delta E$/hit \\
\midrule
01 & intermetallic & $-$0.001 / No & $-$0.001 / No & $-$0.070 / Yes & $-$0.073 / Yes \\
02 & intermetallic & 0.000 / No & $-$0.615 / Yes & $-$0.426 / Yes & $-$0.426 / Yes \\
09 & monometallic & $-$0.173 / N/A & 0.000 / N/A & 0.000 / N/A & $-$0.188 / N/A \\
14 & intermetallic & 0.000 / No & $-$0.372 / Yes & $-$0.160 / Yes & $-$0.126 / Yes \\
19 & intermetallic & $-$0.709 / Yes & $-$0.451 / Yes & 0.000 / No & 0.000 / No \\
\midrule
Support rate & --- & 1/4 = 25\% & 3/4 = 75\% & 3/4 = 75\% & 3/4 = 75\% \\
\bottomrule
\end{tabular}
\end{table}

% Table 7: GPT-5.4 one-shot provenance summary
\begin{table}[t]
\centering
\caption{GPT-5.4 backend provenance summary. Corrected one-shot uses the successful retry for case~08 while retaining case~06 as a reproducible dissociation failure.}
\label{tab:gpt54-openai-summary}
\begin{tabular}{lrrrr}
\toprule
Dataset & Total & Valid & Failed & Mean tokens/run \\
\midrule
Raw one-shot & 20 & 18 & 2 & 4815.2 \\
Retry-only & 2 & 1 & 1 & 5135.0 \\
Corrected one-shot & 20 & 19 & 1 & 4847.6 \\
Ablation matrix & 25 & 25 & 0 & 22198.6 \\
\bottomrule
\end{tabular}
\end{table}

% Table 8: Claude Sonnet 4.6 provenance summary
\begin{table}[t]
\centering
\caption{Claude~Sonnet~4.6 backend provenance summary. Cases~06 and~08 failed due to NNH dissociation (same failure mode as GPT-5.4); no retries were attempted.}
\label{tab:claude-summary}
\begin{tabular}{lrrrr}
\toprule
Dataset & Total & Valid & Failed & Mean tokens/run \\
\midrule
Raw one-shot & 20 & 18 & 2 & 6390.3 \\
Ablation matrix & 25 & 25 & 0 & 29082.0 \\
\bottomrule
\end{tabular}
\end{table}


Tables~\ref{tab:ablation-gemini} and~\ref{tab:ablation-grok} show that the effect of the agentic loop is strongly case- and backend-dependent rather than uniformly monotonic. The Gemini backend is already robust on this subset: the full agent and the 1-shot baseline coincide on cases 02 and 14, and the other ablations remain effectively identical to full on cases 01, 02, and 14. Case 09 still benefits from the loop, with 1-shot trailing full by 0.173~eV. The main Gemini failure mode is case 19, where disabling termination degrades the final energy by 0.247~eV, while removing the \textit{FORBID} text unexpectedly improves the result by 0.201~eV relative to the nominal full configuration. This is an important negative result: explicit negative constraints are not universally beneficial and can over-prune exploration on some reactive surfaces.

The Grok-4 backend exhibits a clearer closed-loop gain. Relative to the full agent, the 1-shot baseline is worse by 0.615~eV on case 02, 0.372~eV on case 14, and 0.451~eV on case 19. On the same hardest case, removing slip-conditioned feedback or removing the \textit{FORBID} instruction also degrades the final non-dissociated solution by about 0.451~eV. In other words, when the base planner is less stable, the memory of failed geometric intents becomes functionally important rather than merely cosmetic.

\subsection{Termination mostly controls compute, not chemistry}

The termination signal behaves primarily as a cost-control mechanism. For Grok-4, disabling termination matches the full agent on four of the five cases and is only 0.017~eV worse on case 19. For Gemini, the same ablation preserves the full energy on four of the five cases and fails noticeably only on case 19. However, the compute overhead is substantial: disabling termination increases token consumption from 26.7k to 44.6k on Grok case 09, from 21.4k to 38.9k on Grok case 19, and from 29.1k to 170.0k on Gemini case 09. Thus, the termination logic should be understood mainly as a guard against planner churn after the search has already converged.

This reading is consistent with the statistical outputs of the completed matrices. On the Grok-4 ablation, the overall Friedman test across variants is significant ($p = 0.0185$), but pairwise Wilcoxon comparisons remain non-significant after correction because the representative set contains only five cases. Gemini shows no overall variant effect on the same test ($p = 0.3006$). We therefore interpret the ablation primarily through effect sizes and reproducible case-local failures, rather than through binary significance labels.

\subsection{Closed-loop search reduces backend variance}

Table~\ref{tab:cross-llm-summary} summarizes agreement between Gemini and Grok on the same 25 case-variant rows. Across the entire matrix, the median absolute energy difference is 0.0~eV, with exact agreement on 15/25 rows and agreement within 0.01~eV on 17/25 rows. The backend variance is smallest for the iterative full agent (mean $|\Delta E| = 0.041$~eV, 4/5 exact matches) and largest for the 1-shot baseline (mean $|\Delta E| = 0.324$~eV, 1/5 exact matches). The largest disagreement in the full comparison does not occur in the iterative setting, but in 1-shot case 02, where Gemini and Grok differ by 0.615~eV. These data support a concrete systems claim: the closed-loop architecture is not only a search strategy, but also a backend-robustness mechanism.

The remaining disagreement is concentrated rather than diffuse. Case 19 is the dominant stress test: it is the only representative system that remains backend-sensitive across all five variants. This concentration is scientifically useful because it tells us where the current agent is brittle. The problem is not generic adsorption search, but a specific combination of a large adsorbate, a reactive alloy surface, and unstable intermediate plans.

\subsection{Scope of external comparison}

We intentionally avoid treating raw energy differences against published Adsorb-Agent values as evidence of superior adsorption strength. The published Adsorb-Agent numbers were obtained under a different MLFF and reference-energy scheme, whereas the results reported here use MACE throughout. Directly subtracting those energies would therefore conflate search behavior with incompatible energetic conventions. In the current manuscript, the main evidence for the agentic system is therefore internal: controlled ablations, backend-robustness analysis, and explicit failure-case tracing under a fixed physical backend.
